V~současném ekonomickém prostředí, charakterizovaném rychlým technologickým vývojem a zvyšující se
mírou nejistoty, mohou sociální dialog a kolektivní vyjednávání hrát klíčovou roli pro úspěšnou adaptaci na
tento vývoj zajištěním efektivní transparentní spolupráce sociálních partnerů. Orientace trhu práce v ČR
v rámci nadnárodních dodavatelsko-odběratelských vztahů na levnou práci představuje zásadní výzvu pro
hospodářskou politiku země v okamžiku, kdy se u takových vztahů objevují nové tlaky na cenu práce.